\section{Extensions and Future Work}
\label{sec:extensions}

The framework of \cref{sec:framework} is deliberately minimal so that
extensions can be plugged in without rewriting the agent skeleton. Five
directions appear especially fruitful and we sketch each below.

\subsection{Bias-corrected clearing}
\label{sec:ext-clearing}

\Cref{sec:disc-balancer} argued that the bid-selection asymmetry
embedded in the max-bid clearing rule cannot be cancelled by any
linear balancer. Two replacements for the clearing functional $A$
in \cref{eq:price-functional} would address it directly:

\begin{itemize}[leftmargin=*,nosep]
\item \textbf{Winsorised mean}: replace $\bar{r}_{t}$ in
      \cref{eq:value-update} with a 10\% winsorised mean of the yellow
      buyers' winning ratios. This caps the influence of extreme
      single-epoch outliers without changing the agent rule set.
\item \textbf{Selection-aware estimator}: substitute the median or a
      lower-percentile statistic of all bids (not only winning bids)
      for $\bar{r}_{t}$. This estimates the underlying value
      distribution rather than the maximum-order statistic.
\end{itemize}

Both modifications are one-line edits to \texttt{EDEMModel.\_end\_epoch}
in our reference implementation. We expect either to break the
exponential bubble of Run~6 entirely while leaving the disequilibrium
regimes of Runs 2--5 qualitatively intact, since those derive from
schedule dynamics rather than from clearing asymmetry.

\subsection{Realtors and commission structures}
\label{sec:ext-realtors}

The DE seller is a strict utility-maximiser whose only friction is
patience. Real-world residential transactions are intermediated by
realtors whose payoff is a fraction of the sale price; this ties
realtor incentives to seller incentives but injects an additional
clock (the listing agreement) and a different bargaining geometry
(realtors negotiate on behalf of multiple sellers in parallel). A
realtor extension would replace the seller agent with an aggregated
listing agent whose patience and ask-price strategies optimise
expected commission across a portfolio. This is a natural setting
for empirical calibration: realtor commission rates and listing
durations are publicly observable, unlike the patience parameter of
the bare DE seller.

\subsection{Open vs.\ blind auctions}
\label{sec:ext-auctions}

EDEM as posed is a blind-bid model: a buyer's bid does not depend on
other buyers' bids on the same home. Open-bid auctions
(``best-and-final'' rounds, escalation clauses) substantially change
the bid distribution by exposing buyers to one another's valuations.
The likely effect is to amplify the bid-selection asymmetry in
\cref{sec:framework-de-special}, because each buyer can revise upward
in response to seeing competitors. Conversely, a sealed-bid auction
with full information disclosure post-clearing would let buyers
calibrate over time and could reduce the asymmetry. Both regimes
are minor modifications to the buyer agent's bid procedure.

\subsection{Construction and the supply side}
\label{sec:ext-construction}

The construction-market variant invertes the EDEM agent roles:
\emph{contractors} bid downward to win projects, with the
lowest-bid contractor winning. The selection asymmetry then operates
in reverse, biasing realised contract prices below the population
fair value, and the dynamics of contractor entry / exit are governed
by an asymmetric balancer (a contractor exits after sustained losses,
a homeowner does not exit a partially-built project lightly). This
is a one-class extension that captures the historically high
bankruptcy rates of small construction companies as a stable
emergent feature rather than as an exogenous shock.

\subsection{Adaptive agents}
\label{sec:ext-adaptive}

EDEM agents are non-learning. Two adaptive extensions are obvious:

\begin{itemize}[leftmargin=*,nosep]
\item \textbf{Bayesian estimation.} Each agent maintains a posterior
      over the home's fair value and updates it after every observed
      sale. This converges, in the limit, toward the bias-corrected
      regime sketched in \cref{sec:ext-clearing}, since the agent
      learns the selection structure of observed prices.
\item \textbf{Reinforcement-learning bidders.} Each agent is a
      policy mapping local state (price level, time on market,
      recent sales) to bids. With a buy-low/sell-high reward
      structure, RL agents would in principle discover the
      bid-selection asymmetry and exploit it; this would be a
      computational analogue to the Zillow-Offers iBuying problem
      \citep{zillow2021ibuying} and would offer a stress test of any
      proposed fix in \cref{sec:ext-clearing}.
\end{itemize}

We have implemented none of these in the present codebase but each
is a self-contained extension on top of the existing
\texttt{Buyer}/\texttt{Seller} classes; pull requests are welcomed at
the project repository.
